\begin{abstract}
Point forecasts trained with mean squared error are systematically smoother than the series they predict, and a long line of work treats this over-smoothing as a defect to be repaired. We ask a prior question: is the observed smoothness always fixable? We decompose the future into a predictable conditional mean and an irreducible noise term, and show that the Bayes-optimal forecast must already look smoother than any realized series. Over-smoothing therefore has two components: rational shrinkage, which is the optimal response to uncertainty, and unexploited structure, which is repairable. We operationalize the distinction with the \emph{oracle gap}, a training-free diagnostic that simulates the signature an oracle would exhibit on a given dataset and compares it with the signature of a trained model. Across 14 dataset-model cells and a 665-cell audit of 13 time-series foundation models on 50 benchmark variants, the oracle gap predicts when anti-smoothing interventions help, with no contradictory cell. We then introduce \rfc, a region-weighted loss that up-weights points where the prediction under-shoots the realized amplitude. Guided by the diagnostic, \rfc~reduces MSE by 2.3\% on Electricity across three seeds and four horizons, and degrades gracefully to the MSE baseline exactly where the gap is non-positive. For regimes where smoothing is rational, we show that a light-weight descriptor head that predicts structural properties of the future window cannot repair point accuracy, but provides a reliable forward-looking signal for selective prediction: abstaining on the 20\% highest-risk windows reduces the remaining MSE by up to 5.3\%. Our results replace a blanket intuition about over-smoothing with a measurable and predictable condition for when it is worth fixing.
\end{abstract}
